\scanpage{144}
% lines 1–6
On obtient $E$ ainsi : \deleted{\unsure{collapsant|collapsons} $Y$} On définit un schéma en courbes plat $Z$ sur $S$ comme conoyau de
\diagram{\begin{tikzcd}[ampersand replacement=\&,column sep=large]
X\amalg(Y\times_S Y)\arrow[r,shift left=1.15ex,"{\mathrm{id}_X,\alpha\operatorname{pr}_1^{Y/S}}"]\arrow[r,shift right=1.15ex,"{(\mathrm{id}_X,\alpha\operatorname{pr}_2^{Y/S})}"'] \& X \& (\alpha:Y\hookrightarrow X)
\end{tikzcd}}
et on définit
\[
E=\underline{\mathrm{Pic}}^{0}_{Z/S}.
\]

% lines 7–11
\deleted{\unsure{Remarque}} Le morphisme $E\to J$ est induit par $X\xrightarrow{\varphi} Z$ en passant au $\underline{\mathrm{Pic}}^{0}_{-/S}$. Notons que nous avons un torseur canonique $P$ sous $E$,
\[
P=\underline{\mathrm{Pic}}^{1}_{Z/S}
\]
(et, d'ailleurs, un $S$-morphisme can $U\to P$).

% lines 12–19
Notons que $P$ définit un élément
\[
\boxed{\unsure{\eta|y}\in H^1(S,E)},
\]
et la relation entre $\xi_n$ et $\eta$ est
\[
\boxed{\xi_n=\delta_n(\eta)},
\]
\margin{à vérifier \unsure{......}}
où $\delta_n$ est l'opérateur cobord
\[
H^1(S,E)\unsure{-|\to}H^{\unsure{2|1}}(S;{}_{n}E)
\]
associé à la suite exacte de Kummer
\[
0\to{}_{n}E\to E\overset{n\,\mathrm{id}_E}{-}E-0.
\]
% archive sheet number: 143

